\section{El sistema masa--resorte}

Uno de los sistemas físicos más sencillos en los que aparece el movimiento armónico simple está formado por una masa unida a un resorte.

Este sistema resulta especialmente importante porque permite conectar ideas ya conocidas:

\begin{itemize}
    \item la fuerza elástica y la ley de Hooke;
    \item la segunda ley de Newton;
    \item el movimiento armónico simple;
    \item el periodo, la frecuencia y la frecuencia angular.
\end{itemize}

La pregunta central será: ¿por qué una masa unida a un resorte puede oscilar?

\begin{idea}[¿Por qué una masa unida a un resorte puede oscilar?]
Cuando un resorte se estira o se comprime respecto de su condición de equilibrio, ejerce una fuerza que tiende a devolver el sistema hacia dicha condición.

Si la masa posee velocidad al alcanzar el equilibrio, su inercia hace que continúe moviéndose y deforme el resorte en el sentido contrario. Aparece entonces una nueva fuerza restauradora que vuelve a acelerar la masa hacia el equilibrio.

Bajo determinadas condiciones, este proceso se repite y constituye un movimiento armónico simple.
\end{idea}


\subsection{Modelo ideal de un sistema masa--resorte}

Comenzaremos con el caso más sencillo: un bloque de masa \(m\) unido a un resorte horizontal.

Supondremos que:

\begin{itemize}
    \item el bloque puede tratarse como una partícula;
    \item el resorte tiene masa despreciable;
    \item el resorte cumple la ley de Hooke;
    \item la superficie es horizontal y sin rozamiento;
    \item el movimiento ocurre únicamente a lo largo del eje \(x\);
    \item no actúan fuerzas externas variables sobre el sistema.
\end{itemize}

\begin{modelo}[Sistema masa--resorte ideal]
Un sistema masa--resorte ideal está formado por una masa \(m\) unida a un resorte de constante elástica \(k\), cuya masa se desprecia y que obedece la ley de Hooke.

Si además se desprecia el rozamiento, el sistema puede realizar movimiento armónico simple.
\end{modelo}


\subsection{Posición de equilibrio}

Cuando el resorte horizontal se encuentra en su longitud natural y no existe ninguna otra fuerza horizontal, el bloque permanece en reposo.

Elegiremos esta posición como origen:

\[ x=0. \]

La coordenada \(x\) medirá el desplazamiento de la masa respecto de la posición de equilibrio.

Tomaremos el sentido positivo del eje \(x\) hacia la derecha.

\begin{center}
\begin{tikzpicture}[scale=1.0]

    \draw[very thick] (-4,-0.8) -- (4,-0.8);

    \draw[very thick] (-4,-0.8) -- (-4,1.2);
    \foreach \y in {-0.6,-0.2,0.2,0.6,1.0}{
        \draw (-4,\y) -- (-4.25,\y-0.18);
    }

    \draw[
        thick,
        decorate,
        decoration={coil,aspect=0.45,segment length=5pt,amplitude=5pt}
    ] (-4,0.2) -- (-0.8,0.2);

    \draw[fill=GrisFondo,thick] (-0.8,-0.35) rectangle (0.6,0.75);
    \node at (-0.1,0.2) {\(m\)};

    \draw[dashed] (-0.1,-1.1) -- (-0.1,1.3);
    \node[below] at (-0.1,-1.1) {\(x=0\)};

    \draw[-{Latex},thick] (-0.1,-1.45) -- (3.5,-1.45);
    \node[right] at (3.5,-1.45) {\(+x\)};

    \node[above] at (-2.3,0.75) {resorte};

\end{tikzpicture}
\end{center}

Si desplazamos la masa hacia la derecha:

\[ x>0, \]

el resorte queda estirado y ejerce una fuerza hacia la izquierda.

Si desplazamos la masa hacia la izquierda:

\[ x<0, \]

el resorte queda comprimido y ejerce una fuerza hacia la derecha.

En ambos casos, la fuerza apunta hacia la posición de equilibrio.


\section{Fuerza restauradora del resorte}

Dentro del régimen elástico lineal, la ley de Hooke establece:

\[ \boxed{F_{s,x}=-kx}. \]

Aquí:

\[
F_{s,x}=\text{componente de la fuerza ejercida por el resorte},
\]

\[
k=\text{constante elástica del resorte},
\]

\[
x=\text{desplazamiento respecto del equilibrio}.
\]

La unidad SI de \(k\) es:

\[ [k]=\si{\newton\per\meter}. \]

El signo negativo contiene información física fundamental.

\begin{interpretacion}[Significado del signo negativo]
En

\[ F_{s,x}=-kx, \]

el signo negativo indica que la fuerza del resorte tiene sentido opuesto al desplazamiento.

Por tanto:

\[
x>0
\quad\Longrightarrow\quad
F_{s,x}<0,
\]

mientras que:

\[
x<0
\quad\Longrightarrow\quad
F_{s,x}>0.
\]

La fuerza siempre tiende a devolver la masa hacia \(x=0\).
\end{interpretacion}

\begin{advertencia}
No debe escribirse siempre

\[ F_s=kx \]

si se está trabajando con componentes y signos.

La expresión

\[ F_s=k|x| \]

puede utilizarse para la \textbf{magnitud} de la fuerza, mientras que

\[ F_{s,x}=-kx \]

expresa correctamente su componente y su sentido sobre el eje elegido.
\end{advertencia}


\section{Segunda ley de Newton aplicada al sistema}

Sobre el bloque horizontal actúan:

\begin{itemize}
    \item el peso \(\vect{P}=m\vect{g}\);
    \item la fuerza normal \(\vect{N}\);
    \item la fuerza elástica \(\vect{F}_s\).
\end{itemize}

En la dirección vertical no existe aceleración, por lo que:

\[ \sum F_y=0. \]

Así:

\[ N-mg=0. \]

Por tanto:

\[ N=mg. \]

Estas fuerzas se equilibran verticalmente y no determinan la oscilación horizontal.

En el eje \(x\), la única fuerza es la fuerza elástica:

\[ \sum F_x=F_{s,x}. \]

Aplicando la segunda ley de Newton:

\[ \sum F_x=ma_x. \]

Entonces:

\[ -kx=ma_x. \]

Despejando la aceleración:

\[ a_x=-\frac{k}{m}x. \]

Esta relación tiene exactamente la forma característica del movimiento armónico simple:

\[ a_x=-\omega^2x. \]

Por comparación:

\[ \omega^2=\frac{k}{m}. \]

Por tanto:

\[ \boxed{\omega=\sqrt{\frac{k}{m}}}. \]

\begin{leyfisica}[Sistema masa--resorte ideal]
Una masa \(m\) unida a un resorte ideal de constante elástica \(k\), sin rozamiento, realiza movimiento armónico simple alrededor de su posición de equilibrio con frecuencia angular

\[ \boxed{\omega=\sqrt{\frac{k}{m}}}. \]
\end{leyfisica}


\subsection{Comprobación dimensional}

La constante elástica posee unidades:

\[ [k]=\frac{\si{\newton}}{\si{\meter}}. \]

Como:

\[ \SI{1}{\newton}=\SI{1}{\kilogram\meter\per\second\squared}, \]

tenemos:

\[
[k]
=
\frac{\si{\kilogram\meter\per\second\squared}}{\si{\meter}}
=
\si{\kilogram\per\second\squared}.
\]

Entonces:

\[
\left[\frac{k}{m}\right]
=
\frac{\si{\kilogram\per\second\squared}}{\si{\kilogram}}
=
\si{\per\second\squared}.
\]

Por tanto:

\[
\left[\sqrt{\frac{k}{m}}\right]
=
\si{\per\second},
\]

que corresponde a las dimensiones de una frecuencia angular.


\section{Periodo y frecuencia del sistema masa--resorte}

Sabemos que:

\[ \omega=\frac{2\pi}{T}. \]

Para el sistema masa--resorte:

\[ \omega=\sqrt{\frac{k}{m}}. \]

Entonces:

\[ \frac{2\pi}{T}=\sqrt{\frac{k}{m}}. \]

Despejando:

\[ \boxed{T=2\pi\sqrt{\frac{m}{k}}}. \]

Como:

\[ f=\frac{1}{T}, \]

obtenemos:

\[ \boxed{f=\frac{1}{2\pi}\sqrt{\frac{k}{m}}}. \]

Así, las tres relaciones fundamentales son:

\[
\boxed{
\omega=\sqrt{\frac{k}{m}},
\qquad
T=2\pi\sqrt{\frac{m}{k}},
\qquad
f=\frac{1}{2\pi}\sqrt{\frac{k}{m}}
}.
\]


\subsection{¿De qué depende el periodo?}

La ecuación

\[ T=2\pi\sqrt{\frac{m}{k}} \]

permite interpretar físicamente el sistema.

Si la masa aumenta:

\[ T\propto\sqrt{m}. \]

Una masa mayor posee mayor inercia y responde más lentamente a la fuerza del resorte.

Si la constante elástica aumenta:

\[ T\propto\frac{1}{\sqrt{k}}. \]

Un resorte más rígido produce una fuerza restauradora mayor para el mismo desplazamiento y, por tanto, el sistema oscila más rápidamente.

\begin{interpretacion}
Para un sistema masa--resorte ideal:

\[
m\uparrow
\quad\Longrightarrow\quad
T\uparrow,\qquad f\downarrow,
\]

mientras que:

\[
k\uparrow
\quad\Longrightarrow\quad
T\downarrow,\qquad f\uparrow.
\]
\end{interpretacion}


\subsection{Dependencia con la amplitud}

Observe que la amplitud \(A\) no aparece en:

\[ T=2\pi\sqrt{\frac{m}{k}}. \]

Por tanto, dentro del modelo ideal:

\begin{principio}[Independencia del periodo respecto de la amplitud]
Mientras el resorte permanezca dentro del régimen donde se cumple la ley de Hooke, el periodo de un sistema masa--resorte ideal no depende de la amplitud.
\end{principio}

Esto significa que dos oscilaciones del mismo sistema con amplitudes diferentes tienen el mismo periodo, siempre que ambas permanezcan dentro del régimen lineal del resorte.

\begin{advertencia}
La independencia respecto de la amplitud es una consecuencia del modelo lineal:

\[ F_{s,x}=-kx. \]

Si la deformación es suficientemente grande para que el resorte deje de obedecer la ley de Hooke, el movimiento deja de ser exactamente armónico simple y el periodo puede depender de la amplitud.
\end{advertencia}


\section{Ecuación del movimiento de una masa unida a un resorte}

Para cualquier movimiento armónico simple:

\[ x(t)=A\cos(\omega t+\phi). \]

En un sistema masa--resorte:

\[ \omega=\sqrt{\frac{k}{m}}. \]

Por tanto:

\[ \boxed{x(t)=A\cos\left(\sqrt{\frac{k}{m}}\,t+\phi\right)}. \]

La velocidad es:

\[ v_x(t)=-\omega A\sin(\omega t+\phi). \]

Entonces:

\[ \boxed{v_x(t)=-A\sqrt{\frac{k}{m}}\sin\left(\sqrt{\frac{k}{m}}\,t+\phi\right)}. \]

La aceleración es:

\[ a_x=-\omega^2x. \]

Como:

\[ \omega^2=\frac{k}{m}, \]

obtenemos:

\[ \boxed{a_x=-\frac{k}{m}x}. \]

O, explícitamente en función del tiempo:

\[ \boxed{a_x(t)=-\frac{k}{m}A\cos\left(\sqrt{\frac{k}{m}}\,t+\phi\right)}. \]


\subsection{Caso especial: liberación desde un extremo}

Una situación muy frecuente ocurre cuando la masa se lleva hasta:

\[ x=A \]

y se libera desde el reposo en:

\[ t=0. \]

Entonces:

\[ x(0)=A \]

y:

\[ v_x(0)=0. \]

Estas condiciones se satisfacen con:

\[ \phi=0. \]

Por tanto:

\[ \boxed{x(t)=A\cos(\omega t)}. \]

Para el sistema masa--resorte:

\[ \boxed{x(t)=A\cos\left(\sqrt{\frac{k}{m}}\,t\right)}. \]

La velocidad será:

\[ \boxed{v_x(t)=-\omega A\sin(\omega t)}, \]

y la aceleración:

\[ \boxed{a_x(t)=-\omega^2A\cos(\omega t)}. \]


\section{Interpretación física de una oscilación completa}

Supongamos que la masa se desplaza inicialmente hasta:

\[ x=+A \]

y se libera desde el reposo.

En ese instante:

\[ v_x=0. \]

La fuerza elástica es:

\[ F_{s,x}=-kA, \]

por lo que la aceleración apunta hacia el equilibrio.

La masa comienza a moverse hacia la izquierda.

Al acercarse a:

\[ x=0, \]

la deformación disminuye y también disminuye la magnitud de la fuerza restauradora.

Cuando la masa atraviesa el equilibrio:

\[ x=0, \]

tenemos:

\[ F_{s,x}=0 \]

y:

\[ a_x=0. \]

Sin embargo, la masa no se detiene. Debido a su inercia, continúa moviéndose hacia la izquierda.

El resorte comienza entonces a comprimirse.

Cuando alcanza:

\[ x=-A, \]

la velocidad vuelve a ser cero, pero la fuerza y la aceleración poseen magnitud máxima y apuntan nuevamente hacia el equilibrio.

La masa invierte su movimiento y el proceso se repite.

\begin{center}
\begin{tabularx}{0.94\textwidth}{cYYY}
\toprule
\textbf{Posición} & \textbf{Fuerza} & \textbf{Velocidad} & \textbf{Aceleración} \\
\midrule
\(x=+A\) & máxima hacia \(-x\) & \(0\) & máxima hacia \(-x\) \\
\(x=0\) & \(0\) & máxima en magnitud & \(0\) \\
\(x=-A\) & máxima hacia \(+x\) & \(0\) & máxima hacia \(+x\) \\
\bottomrule
\end{tabularx}
\end{center}

\begin{errorfrecuente}
En la posición de equilibrio la fuerza neta y la aceleración son cero, pero la masa no está necesariamente en reposo.

Durante una oscilación, la rapidez es máxima cuando la masa atraviesa el equilibrio.
\end{errorfrecuente}


\section{Ejemplo: sistema horizontal}

\begin{ejemplo}[Oscilación horizontal de un bloque]
Un bloque de masa

\[ m=\SI{0.500}{\kilogram} \]

está unido a un resorte horizontal de constante elástica

\[ k=\SI{200}{\newton\per\meter}. \]

La superficie es horizontal y sin rozamiento. El bloque se desplaza

\[ A=\SI{8.0}{\centi\meter} \]

hacia la derecha y se libera desde el reposo.

Determinemos la frecuencia angular, el periodo, la frecuencia y la ecuación de posición.
\end{ejemplo}

\subsubsection*{Datos}

Primero expresamos la amplitud en unidades SI:

\[ A=\SI{0.080}{\meter}. \]

Tenemos:

\[
m=\SI{0.500}{\kilogram},
\qquad
k=\SI{200}{\newton\per\meter}.
\]

\subsubsection*{Modelo físico}

El bloque está unido a un resorte ideal y no existe rozamiento significativo.

Por tanto, el movimiento es armónico simple y:

\[ \omega=\sqrt{\frac{k}{m}}. \]

\subsubsection*{Frecuencia angular}

Sustituyendo:

\[
\omega
=
\sqrt{
\frac{\SI{200}{\newton\per\meter}}
{\SI{0.500}{\kilogram}}
}.
\]

Entonces:

\[ \boxed{\omega=\SI{20.0}{\radian\per\second}}. \]

\subsubsection*{Periodo}

\[ T=\frac{2\pi}{\omega}. \]

Por tanto:

\[
T
=
\frac{2\pi}
{\SI{20.0}{\radian\per\second}}
\approx
\SI{0.314}{\second}.
\]

Así:

\[ \boxed{T\approx\SI{0.314}{\second}}. \]

\subsubsection*{Frecuencia}

\[ f=\frac{1}{T}. \]

Entonces:

\[ \boxed{f\approx\SI{3.18}{\hertz}}. \]

\subsubsection*{Ecuación del movimiento}

Como el bloque comienza en:

\[ x=A \]

y se libera desde el reposo:

\[ \phi=0. \]

Por tanto:

\[ x(t)=A\cos(\omega t). \]

Sustituyendo:

\[ \boxed{x(t)=(\SI{0.080}{\meter})\cos(20.0t)}. \]

La velocidad es:

\[ v_x(t)=-\omega A\sin(\omega t), \]

por lo que:

\[ \boxed{v_x(t)=-(\SI{1.60}{\meter\per\second})\sin(20.0t)}. \]

La aceleración es:

\[ a_x(t)=-\omega^2A\cos(\omega t), \]

de modo que:

\[ \boxed{a_x(t)=-(\SI{32.0}{\meter\per\second\squared})\cos(20.0t)}. \]

\begin{interpretacion}
El sistema completa un ciclo aproximadamente cada \(\SI{0.314}{\second}\).

En los extremos, la masa se detiene momentáneamente y su aceleración tiene magnitud máxima.

Al atravesar el equilibrio, su aceleración es cero y su rapidez es máxima.
\end{interpretacion}


\section{Efecto de cambiar la masa o el resorte}

Las relaciones:

\[ \omega=\sqrt{\frac{k}{m}} \]

y:

\[ T=2\pi\sqrt{\frac{m}{k}} \]

permiten realizar predicciones sin efectuar cálculos completos.


\subsection{Cambio de masa}

Supongamos que se mantiene el mismo resorte y la masa cambia de \(m\) a \(4m\).

El nuevo periodo es:

\[
T'
=
2\pi\sqrt{\frac{4m}{k}}.
\]

Entonces:

\[ T'=2T. \]

La frecuencia se reduce a la mitad:

\[ f'=\frac{f}{2}. \]

\begin{interpretacion}
Cuadruplicar la masa duplica el periodo.

La dependencia no es lineal, sino proporcional a \(\sqrt{m}\).
\end{interpretacion}


\subsection{Cambio de constante elástica}

Si se mantiene la masa y se cambia el resorte por otro con:

\[ k'=4k, \]

entonces:

\[
T'
=
2\pi\sqrt{\frac{m}{4k}}
=
\frac{T}{2}.
\]

Por tanto:

\[ f'=2f. \]

Un resorte cuatro veces más rígido produce una frecuencia dos veces mayor.

\begin{ejemplo}[Comparación sin cálculo completo]
Dos sistemas utilizan resortes idénticos. El primero contiene una masa \(m\) y el segundo una masa \(9m\).

Como:

\[ T\propto\sqrt{m}, \]

tenemos:

\[
\frac{T_2}{T_1}
=
\sqrt{\frac{9m}{m}}
=
3.
\]

Por tanto:

\[ \boxed{T_2=3T_1}. \]

La frecuencia del segundo sistema será:

\[ \boxed{f_2=\frac{f_1}{3}}. \]

No fue necesario conocer el valor de \(k\).
\end{ejemplo}


\section{Sistema masa--resorte vertical}

El sistema también puede disponerse verticalmente.

Ahora la gravedad actúa a lo largo de la dirección de movimiento, por lo que debemos analizar cuidadosamente la posición de equilibrio.


\subsection{Longitud natural y posición de equilibrio}

Consideremos un resorte vertical de constante elástica \(k\).

Antes de colocar la masa, el resorte posee una longitud natural.

Al colgar una masa \(m\), el resorte se estira hasta alcanzar una nueva posición de equilibrio.

Representaremos mediante:

\[ \Delta\ell_{\mathrm{eq}} \]

el alargamiento estático producido por el peso.

\begin{center}
\begin{tikzpicture}[scale=0.95]

    \draw[very thick] (-2.0,2.5) -- (2.0,2.5);

    \foreach \x in {-1.8,-1.4,...,1.8}{
        \draw (\x,2.5) -- (\x-0.25,2.75);
    }

    \draw[
        thick,
        decorate,
        decoration={coil,aspect=0.45,segment length=5pt,amplitude=5pt}
    ] (0,2.5) -- (0,-0.2);

    \draw[fill=GrisFondo,thick] (-0.65,-0.2) rectangle (0.65,-1.2);
    \node at (0,-0.7) {\(m\)};

    \draw[-{Latex},thick] (1.8,1.3) -- (1.8,-1.8);
    \node[right] at (1.8,-1.8) {\(+y\)};

    \draw[dashed] (-1.3,-0.7) -- (1.3,-0.7);
    \node[left] at (-1.3,-0.7) {equilibrio};

\end{tikzpicture}
\end{center}

Elegiremos ahora el eje \(y\) positivo hacia abajo.

En equilibrio:

\[ a_y=0. \]

Las fuerzas sobre la masa son:

\begin{itemize}
    \item el peso \(mg\), hacia abajo;
    \item la fuerza del resorte, hacia arriba.
\end{itemize}

Por tanto:

\[ \sum F_y=0. \]

Así:

\[ mg-k\Delta\ell_{\mathrm{eq}}=0. \]

Entonces:

\[ \boxed{k\Delta\ell_{\mathrm{eq}}=mg}. \]

De aquí:

\[ \boxed{\Delta\ell_{\mathrm{eq}}=\frac{mg}{k}}. \]

\begin{interpretacion}
En un sistema vertical, el peso modifica la posición de equilibrio.

Cuanto mayor sea la masa, mayor será el alargamiento estático.

Cuanto más rígido sea el resorte, menor será dicho alargamiento.
\end{interpretacion}


\subsection{Oscilaciones alrededor del nuevo equilibrio}

Ahora desplazemos la masa una distancia \(y\) respecto de su posición de equilibrio.

Elegimos:

\[ y=0 \]

en el equilibrio y \(+y\) hacia abajo.

La deformación total del resorte será:

\[ \Delta\ell=\Delta\ell_{\mathrm{eq}}+y. \]

La segunda ley de Newton da:

\[ ma_y=mg-k(\Delta\ell_{\mathrm{eq}}+y). \]

Distribuyendo:

\[ ma_y=mg-k\Delta\ell_{\mathrm{eq}}-ky. \]

Pero en equilibrio:

\[ mg=k\Delta\ell_{\mathrm{eq}}. \]

Por tanto, los dos primeros términos se cancelan:

\[ ma_y=-ky. \]

Entonces:

\[ a_y=-\frac{k}{m}y. \]

Esta es nuevamente la ecuación de un movimiento armónico simple.

Por tanto:

\[ \boxed{\omega=\sqrt{\frac{k}{m}}}, \]

\[ \boxed{T=2\pi\sqrt{\frac{m}{k}}}, \]

y:

\[ \boxed{f=\frac{1}{2\pi}\sqrt{\frac{k}{m}}}. \]

\begin{principio}[Efecto de la gravedad en el sistema vertical]
La gravedad cambia la posición de equilibrio del sistema masa--resorte vertical, pero no cambia su frecuencia angular ni su periodo dentro del modelo ideal.

Una vez que el desplazamiento se mide respecto de la nueva posición de equilibrio, la ecuación del movimiento es:

\[ a_y=-\frac{k}{m}y. \]
\end{principio}

\begin{errorfrecuente}
En un sistema vertical no debe tomarse necesariamente como \(y=0\) la posición donde el resorte está sin deformar.

La oscilación ocurre alrededor de la posición de equilibrio definida por:

\[ k\Delta\ell_{\mathrm{eq}}=mg. \]
\end{errorfrecuente}


\subsection{Periodo a partir del alargamiento estático}

Como:

\[ k\Delta\ell_{\mathrm{eq}}=mg, \]

podemos despejar:

\[ k=\frac{mg}{\Delta\ell_{\mathrm{eq}}}. \]

El periodo es:

\[ T=2\pi\sqrt{\frac{m}{k}}. \]

Sustituyendo:

\[
T
=
2\pi
\sqrt{
\frac{m}
{mg/\Delta\ell_{\mathrm{eq}}}
}.
\]

Simplificando:

\[ \boxed{T=2\pi\sqrt{\frac{\Delta\ell_{\mathrm{eq}}}{g}}}. \]

Esta expresión permite determinar el periodo de un sistema vertical si conocemos cuánto se alarga estáticamente el resorte al colocar la masa.

\begin{advertencia}
La expresión

\[ T=2\pi\sqrt{\frac{\Delta\ell_{\mathrm{eq}}}{g}} \]

no significa que el periodo fundamental del sistema dependa directamente de \(g\).

Como:

\[ \Delta\ell_{\mathrm{eq}}=\frac{mg}{k}, \]

el cociente

\[ \frac{\Delta\ell_{\mathrm{eq}}}{g}=\frac{m}{k}. \]

Para un mismo sistema ideal con \(m\) y \(k\) fijos:

\[ T=2\pi\sqrt{\frac{m}{k}}. \]
\end{advertencia}


\section{Ejemplo: resorte vertical}

\begin{ejemplo}[Equilibrio y oscilación vertical]
Una masa de

\[ m=\SI{0.400}{\kilogram} \]

se cuelga de un resorte vertical de constante

\[ k=\SI{100}{\newton\per\meter}. \]

Tomemos:

\[ g=\SI{9.81}{\meter\per\second\squared}. \]

Determinemos:

\begin{enumerate}[label=\alph*)]
    \item el alargamiento de equilibrio;
    \item la frecuencia angular;
    \item el periodo;
    \item la frecuencia.
\end{enumerate}
\end{ejemplo}

\subsubsection*{Alargamiento de equilibrio}

En equilibrio:

\[ k\Delta\ell_{\mathrm{eq}}=mg. \]

Despejando:

\[ \Delta\ell_{\mathrm{eq}}=\frac{mg}{k}. \]

Sustituyendo:

\[
\Delta\ell_{\mathrm{eq}}
=
\frac{
(\SI{0.400}{\kilogram})
(\SI{9.81}{\meter\per\second\squared})
}{
\SI{100}{\newton\per\meter}
}.
\]

Por tanto:

\[ \boxed{\Delta\ell_{\mathrm{eq}}=\SI{0.0392}{\meter}\approx\SI{3.92}{\centi\meter}}. \]

\subsubsection*{Frecuencia angular}

\[ \omega=\sqrt{\frac{k}{m}}. \]

Entonces:

\[
\omega
=
\sqrt{
\frac{100}{0.400}
}\ \si{\radian\per\second}.
\]

Así:

\[ \boxed{\omega\approx\SI{15.8}{\radian\per\second}}. \]

\subsubsection*{Periodo}

\[ T=\frac{2\pi}{\omega}. \]

Por tanto:

\[ \boxed{T\approx\SI{0.397}{\second}}. \]

\subsubsection*{Frecuencia}

\[ f=\frac{1}{T}. \]

Así:

\[ \boxed{f\approx\SI{2.52}{\hertz}}. \]

\begin{interpretacion}
El peso produce un alargamiento permanente de aproximadamente \(\SI{3.92}{\centi\meter}\), estableciendo una nueva posición de equilibrio.

Si la masa se perturba alrededor de esa posición, oscila con un periodo aproximado de \(\SI{0.397}{\second}\).
\end{interpretacion}


\section{Ecuación de movimiento en el sistema vertical}

Si utilizamos \(y\) para medir el desplazamiento respecto de la posición de equilibrio, la ecuación es:

\[ y(t)=A\cos(\omega t+\phi). \]

Como:

\[ \omega=\sqrt{\frac{k}{m}}, \]

tenemos:

\[ \boxed{y(t)=A\cos\left(\sqrt{\frac{k}{m}}\,t+\phi\right)}. \]

Es importante distinguir entre dos cantidades:

\[
\Delta\ell_{\mathrm{eq}}
=
\text{alargamiento estático respecto de la longitud natural},
\]

y:

\[
y
=
\text{desplazamiento instantáneo respecto del equilibrio}.
\]

La deformación total del resorte será:

\[ \boxed{\Delta\ell(t)=\Delta\ell_{\mathrm{eq}}+y(t)}. \]

\begin{errorfrecuente}
No confundir

\[ \Delta\ell_{\mathrm{eq}} \]

con la amplitud \(A\).

El primero representa cuánto se deforma el resorte debido al peso cuando el sistema está en equilibrio.

La amplitud representa el máximo desplazamiento adicional respecto de esa posición de equilibrio.
\end{errorfrecuente}


\subsection{Ejemplo de ecuación vertical}

Supongamos que el sistema anterior se desplaza:

\[ A=\SI{3.0}{\centi\meter} \]

hacia abajo desde el equilibrio y se libera desde el reposo.

Tomando \(+y\) hacia abajo:

\[ y(0)=+A \]

y:

\[ v_y(0)=0. \]

Entonces:

\[ \phi=0. \]

Como:

\[ \omega\approx\SI{15.8}{\radian\per\second}, \]

la ecuación del movimiento es:

\[ \boxed{y(t)=(\SI{0.030}{\meter})\cos(15.8t)}. \]

La deformación total del resorte en cualquier instante será:

\[ \Delta\ell(t)=\Delta\ell_{\mathrm{eq}}+y(t). \]

Por tanto:

\[ \boxed{\Delta\ell(t)=\SI{0.0392}{\meter}+(\SI{0.030}{\meter})\cos(15.8t)}. \]

Observe que la masa oscila alrededor de una posición en la que el resorte ya está estirado.


\section{Determinación experimental de la constante elástica}

La constante \(k\) puede determinarse de varias maneras.


\subsection{Método estático}

Aplicamos fuerzas conocidas y medimos la deformación correspondiente.

Dentro del régimen de Hooke:

\[ F=k\Delta\ell. \]

Por tanto:

\[ \boxed{k=\frac{F}{\Delta\ell}}. \]

Si se representan experimentalmente los valores de \(F\) frente a \(\Delta\ell\), la pendiente de la gráfica corresponde a \(k\).

\begin{center}
\begin{tikzpicture}

\begin{axis}[
    width=0.72\textwidth,
    height=6cm,
    axis lines=left,
    xlabel={\(\Delta\ell\)},
    ylabel={\(F\)},
    xmin=0,
    xmax=5,
    ymin=0,
    ymax=5,
    xtick=\empty,
    ytick=\empty,
    samples=2
]

\addplot[very thick,AzulPrincipal,domain=0:4.5] {x};

\node at (axis cs:3.2,3.8) {\(F=k\Delta\ell\)};

\end{axis}

\end{tikzpicture}
\end{center}

En un resorte ideal, la gráfica es lineal en el intervalo donde se cumple la ley de Hooke.


\subsection{Método dinámico}

También podemos utilizar el periodo de oscilación.

Partimos de:

\[ T=2\pi\sqrt{\frac{m}{k}}. \]

Elevando al cuadrado:

\[ T^2=4\pi^2\frac{m}{k}. \]

Por tanto:

\[ \boxed{T^2=\frac{4\pi^2}{k}m}. \]

Esta expresión tiene la forma de una relación lineal entre \(T^2\) y \(m\).

Si construimos una gráfica:

\[ T^2\quad\text{frente a}\quad m, \]

la pendiente es:

\[ b=\frac{4\pi^2}{k}. \]

Por tanto:

\[ \boxed{k=\frac{4\pi^2}{b}}. \]

\begin{interpretacion}
El método dinámico permite obtener \(k\) midiendo tiempos en lugar de fuerzas.

Para reducir la incertidumbre en la medición del periodo suele ser conveniente medir el tiempo correspondiente a varias oscilaciones completas y dividir entre el número de ciclos.
\end{interpretacion}

\begin{ejemplo}[Determinación de \(k\) mediante el periodo]
Una masa de

\[ m=\SI{0.250}{\kilogram} \]

oscila unida a un resorte y se mide un periodo:

\[ T=\SI{0.500}{\second}. \]

De:

\[ T=2\pi\sqrt{\frac{m}{k}}, \]

elevamos al cuadrado:

\[ T^2=4\pi^2\frac{m}{k}. \]

Despejando:

\[ k=\frac{4\pi^2m}{T^2}. \]

Sustituyendo:

\[
k
=
\frac{
4\pi^2(\SI{0.250}{\kilogram})
}{
(\SI{0.500}{\second})^2
}.
\]

Por tanto:

\[ \boxed{k\approx\SI{39.5}{\newton\per\meter}}. \]
\end{ejemplo}


\section{Resortes equivalentes}

En algunos sistemas una masa puede estar conectada a más de un resorte.

Si la fuerza neta continúa siendo proporcional al desplazamiento:

\[ F_{\text{net},x}=-k_{\mathrm{eq}}x, \]

el sistema seguirá comportándose como un oscilador armónico simple con:

\[ \boxed{\omega=\sqrt{\frac{k_{\mathrm{eq}}}{m}}}. \]

La cantidad

\[ k_{\mathrm{eq}} \]

se denomina \textbf{constante elástica equivalente}.


\subsection{Resortes en paralelo}

Consideremos dos resortes de constantes:

\[ k_1 \]

y:

\[ k_2 \]

que experimentan el mismo desplazamiento \(x\).

Cada uno ejerce una fuerza restauradora:

\[ F_1=-k_1x, \]

\[ F_2=-k_2x. \]

La fuerza total es:

\[ F_{\text{net}}=F_1+F_2. \]

Entonces:

\[ F_{\text{net}}=-(k_1+k_2)x. \]

Comparando con:

\[ F_{\text{net}}=-k_{\mathrm{eq}}x, \]

obtenemos:

\[ \boxed{k_{\mathrm{eq}}=k_1+k_2}. \]

Para \(N\) resortes en paralelo:

\[ \boxed{k_{\mathrm{eq}}=\sum_{i=1}^{N}k_i}. \]

\begin{interpretacion}
Los resortes en paralelo producen un sistema más rígido.

Por tanto:

\[ k_{\mathrm{eq}}>k_1 \]

y:

\[ k_{\mathrm{eq}}>k_2 \]

para dos resortes con constantes positivas.

Para una misma masa, el periodo disminuye.
\end{interpretacion}


\subsection{Una masa entre dos resortes}

Un caso frecuente consiste en colocar una masa entre dos resortes.

Si la masa se desplaza una distancia \(x\) hacia la derecha, un resorte se estira y el otro se comprime.

Ambos ejercen fuerzas restauradoras hacia la izquierda.

La fuerza neta es:

\[ F_{\text{net},x}=-(k_1+k_2)x. \]

Por tanto:

\[ \boxed{k_{\mathrm{eq}}=k_1+k_2}. \]

Y:

\[ \boxed{T=2\pi\sqrt{\frac{m}{k_1+k_2}}}. \]


\subsection{Resortes en serie}

Consideremos ahora dos resortes conectados uno después del otro.

Si la fuerza transmitida es \(F\), entonces:

\[ F=k_1x_1 \]

y:

\[ F=k_2x_2. \]

Las deformaciones son:

\[ x_1=\frac{F}{k_1}, \]

\[ x_2=\frac{F}{k_2}. \]

La deformación total es:

\[ x=x_1+x_2. \]

Por tanto:

\[
x
=
\frac{F}{k_1}
+
\frac{F}{k_2}.
\]

Factorizando:

\[
x
=
F
\left(
\frac{1}{k_1}
+
\frac{1}{k_2}
\right).
\]

Para el resorte equivalente:

\[ F=k_{\mathrm{eq}}x. \]

Entonces:

\[ x=\frac{F}{k_{\mathrm{eq}}}. \]

Comparando:

\[ \boxed{\frac{1}{k_{\mathrm{eq}}}=\frac{1}{k_1}+\frac{1}{k_2}}. \]

Para dos resortes:

\[ \boxed{k_{\mathrm{eq}}=\frac{k_1k_2}{k_1+k_2}}. \]

Para \(N\) resortes en serie:

\[ \boxed{\frac{1}{k_{\mathrm{eq}}}=\sum_{i=1}^{N}\frac{1}{k_i}}. \]

\begin{interpretacion}
Una combinación en serie es menos rígida que cualquiera de los resortes individuales.

Por tanto:

\[ k_{\mathrm{eq}}<k_1 \]

y:

\[ k_{\mathrm{eq}}<k_2. \]

Para una misma masa, una constante equivalente menor produce un periodo mayor.
\end{interpretacion}


\section{Ejemplo: comparación de dos configuraciones}

\begin{ejemplo}[Resortes en paralelo y en serie]
Una masa

\[ m=\SI{0.600}{\kilogram} \]

se conecta a dos resortes:

\[
k_1=\SI{120}{\newton\per\meter},
\qquad
k_2=\SI{180}{\newton\per\meter}.
\]

Comparemos el periodo cuando los resortes se conectan en paralelo y cuando se conectan en serie.
\end{ejemplo}

\subsubsection*{Resortes en paralelo}

La constante equivalente es:

\[ k_{\mathrm{eq}}=k_1+k_2. \]

Entonces:

\[
k_{\mathrm{eq}}
=
120+180
=
\SI{300}{\newton\per\meter}.
\]

El periodo es:

\[
T
=
2\pi
\sqrt{
\frac{0.600}{300}
}.
\]

Por tanto:

\[ \boxed{T_{\mathrm{par}}\approx\SI{0.281}{\second}}. \]


\subsubsection*{Resortes en serie}

La constante equivalente es:

\[ k_{\mathrm{eq}}=\frac{k_1k_2}{k_1+k_2}. \]

Entonces:

\[
k_{\mathrm{eq}}
=
\frac{(120)(180)}{120+180}
=
\SI{72.0}{\newton\per\meter}.
\]

El periodo es:

\[
T
=
2\pi
\sqrt{
\frac{0.600}{72.0}
}.
\]

Por tanto:

\[ \boxed{T_{\mathrm{serie}}\approx\SI{0.574}{\second}}. \]

\begin{interpretacion}
La combinación en paralelo es más rígida y oscila más rápidamente.

La combinación en serie es menos rígida y presenta un periodo mayor.
\end{interpretacion}


\section{Limitaciones del modelo ideal}

La expresión

\[ T=2\pi\sqrt{\frac{m}{k}} \]

se obtiene bajo determinadas hipótesis.

Entre ellas:

\begin{itemize}
    \item el resorte cumple la ley de Hooke;
    \item la masa del resorte es despreciable frente a la masa oscilante;
    \item no existe rozamiento significativo;
    \item no existe resistencia del aire significativa;
    \item las propiedades del resorte permanecen constantes;
    \item el movimiento es esencialmente unidimensional.
\end{itemize}


\subsection{Masa del resorte}

Un resorte real posee masa.

Cuando su masa no es despreciable, distintas partes del resorte se mueven con velocidades diferentes durante la oscilación.

Por ello, no toda la masa del resorte puede tratarse simplemente como si estuviera concentrada en el bloque.

\begin{aproximacion}[Resorte ligero]
La expresión

\[ T=2\pi\sqrt{\frac{m}{k}} \]

supone que la masa del resorte es despreciable frente a la masa unida a él.

Si esta condición no se cumple, debe utilizarse un modelo más preciso que incluya la contribución dinámica del propio resorte.
\end{aproximacion}


\subsection{Deformaciones grandes}

La ley

\[ F=-kx \]

es una aproximación lineal.

Un resorte real puede dejar de comportarse linealmente si se deforma demasiado.

En ese caso:

\[ F\not\propto -x \]

y el movimiento deja de ser exactamente armónico simple.


\subsection{Rozamiento}

Un sistema real presenta normalmente alguna forma de rozamiento o resistencia.

En tal situación, las oscilaciones no continúan indefinidamente con amplitud constante.

Por tanto, el sistema masa--resorte ideal representa una aproximación a sistemas reales durante aquellos intervalos en los que las pérdidas puedan despreciarse.


\section{Aplicaciones del sistema masa--resorte}

El modelo masa--resorte aparece en numerosas situaciones.

Entre ellas:

\begin{itemize}
    \item suspensiones de vehículos;
    \item básculas mecánicas;
    \item vibraciones de máquinas;
    \item sensores;
    \item instrumentos de medición;
    \item modelos de vibraciones moleculares;
    \item estructuras elásticas;
    \item sistemas de aislamiento de vibraciones.
\end{itemize}

En muchos casos el ``resorte'' no es literalmente un resorte helicoidal. Cualquier sistema cuya fuerza restauradora pueda aproximarse mediante:

\[ F=-k_{\mathrm{ef}}x \]

puede comportarse dinámicamente como un sistema masa--resorte.

\begin{aplicacion}[Suspensión de un vehículo]
La suspensión de un automóvil puede modelarse, en una primera aproximación, como una masa apoyada sobre resortes.

Si la masa efectiva aumenta, el periodo de oscilación aumenta.

Si la suspensión utiliza resortes más rígidos, la frecuencia natural aumenta.

El modelo ideal no describe todos los detalles de una suspensión real, pero permite comprender su comportamiento básico.
\end{aplicacion}


\section{Errores frecuentes}

\begin{errorfrecuente}
\textbf{Utilizar la longitud natural como origen en un sistema vertical.}

La masa oscila alrededor de la posición de equilibrio, no alrededor de la posición donde el resorte está sin deformar.
\end{errorfrecuente}

\begin{errorfrecuente}
\textbf{Olvidar el signo de la fuerza elástica.}

La relación por componentes es:

\[ F_{s,x}=-kx. \]

El signo negativo indica que la fuerza es restauradora.
\end{errorfrecuente}

\begin{errorfrecuente}
\textbf{Creer que un resorte más rígido produce un periodo mayor.}

En realidad:

\[ T=2\pi\sqrt{\frac{m}{k}}. \]

Al aumentar \(k\), el periodo disminuye.
\end{errorfrecuente}

\begin{errorfrecuente}
\textbf{Creer que una masa mayor oscila más rápidamente.}

Para un mismo resorte:

\[ T\propto\sqrt{m}. \]

Por tanto, una masa mayor produce un periodo mayor y una frecuencia menor.
\end{errorfrecuente}

\begin{errorfrecuente}
\textbf{Introducir \(g\) innecesariamente en el periodo de un sistema vertical.}

Para un sistema masa--resorte ideal:

\[ T=2\pi\sqrt{\frac{m}{k}}. \]

La gravedad modifica la posición de equilibrio, pero no cambia esta expresión.
\end{errorfrecuente}

\begin{errorfrecuente}
\textbf{Confundir alargamiento de equilibrio y amplitud.}

En un resorte vertical:

\[ \Delta\ell_{\mathrm{eq}}=\frac{mg}{k} \]

describe la deformación estática producida por el peso.

La amplitud \(A\) describe el máximo desplazamiento respecto de esa posición de equilibrio.
\end{errorfrecuente}

\begin{errorfrecuente}
\textbf{Aplicar la fórmula del periodo sin comprobar el modelo.}

La expresión

\[ T=2\pi\sqrt{\frac{m}{k}} \]

requiere un comportamiento aproximadamente lineal, pérdidas despreciables y un resorte cuya masa pueda despreciarse o tratarse mediante una aproximación adecuada.
\end{errorfrecuente}


\section{Secuencia sugerida para impartir la clase}

\begin{notapedagogica}
Una posible secuencia para desarrollar este tema en clase es:

\begin{enumerate}
    \item Mostrar físicamente un resorte y una masa o representar el sistema horizontal.

    \item Preguntar qué ocurre si la masa se desplaza hacia la derecha y se libera.

    \item Definir la posición de equilibrio como \(x=0\).

    \item Recordar brevemente la ley de Hooke:
    \[ F_{s,x}=-kx. \]

    \item Dibujar el diagrama de cuerpo libre y mostrar que peso y normal se equilibran verticalmente.

    \item Aplicar la segunda ley de Newton en el eje \(x\):
    \[ -kx=ma_x. \]

    \item Obtener:
    \[ a_x=-\frac{k}{m}x \]
    y compararla con:
    \[ a_x=-\omega^2x. \]

    \item Deducir:
    \[ \omega=\sqrt{\frac{k}{m}}. \]

    \item Construir a partir de ella:
    \[ T=2\pi\sqrt{\frac{m}{k}} \]
    y:
    \[ f=\frac{1}{2\pi}\sqrt{\frac{k}{m}}. \]

    \item Analizar cualitativamente cómo cambian \(T\) y \(f\) al modificar \(m\) o \(k\).

    \item Resolver un ejemplo numérico horizontal.

    \item Pasar al resorte vertical y preguntar si la gravedad cambia el periodo.

    \item Determinar primero la nueva posición de equilibrio:
    \[ k\Delta\ell_{\mathrm{eq}}=mg. \]

    \item Medir el desplazamiento desde esa nueva posición y demostrar que:
    \[ ma_y=-ky. \]

    \item Resolver un ejemplo vertical.

    \item Mostrar cómo determinar experimentalmente \(k\) mediante una gráfica \(F-\Delta\ell\) o mediante el periodo.

    \item Si el tiempo de clase lo permite, estudiar resortes equivalentes en paralelo y en serie.

    \item Finalizar con preguntas conceptuales que permitan distinguir equilibrio, amplitud, rigidez, masa y periodo.
\end{enumerate}
\end{notapedagogica}


\section{Preguntas de comprensión}

\begin{enumerate}
    \item ¿Por qué la fuerza ejercida por un resorte se denomina fuerza restauradora?

    \item ¿Qué representa físicamente el signo negativo en
    \[ F_{s,x}=-kx? \]

    \item ¿Por qué un bloque sobre una superficie horizontal sin rozamiento realiza MAS cuando está unido a un resorte ideal?

    \item ¿Qué relación existe entre \(k\), \(m\) y \(\omega\)?

    \item Si se duplica la masa y se mantiene el mismo resorte, ¿el periodo se duplica? Justifique.

    \item Si la constante elástica aumenta, ¿la frecuencia aumenta o disminuye?

    \item ¿Por qué el periodo ideal no depende de la amplitud?

    \item ¿Bajo qué condición podría comenzar a depender de la amplitud?

    \item En un sistema horizontal, ¿qué fuerzas actúan verticalmente sobre la masa?

    \item ¿Por qué esas fuerzas no determinan la oscilación horizontal?

    \item ¿En qué punto de la oscilación la fuerza del resorte es cero?

    \item ¿En qué punto la magnitud de la fuerza restauradora es máxima?

    \item En un sistema vertical, ¿qué diferencia existe entre la longitud natural y la posición de equilibrio?

    \item ¿Qué representa
    \[ \Delta\ell_{\mathrm{eq}}=\frac{mg}{k}? \]

    \item ¿Por qué la gravedad modifica la posición de equilibrio pero no el periodo ideal?

    \item ¿Cuál es la diferencia entre \(\Delta\ell_{\mathrm{eq}}\) y la amplitud?

    \item Si dos resortes se colocan en paralelo, ¿el sistema resulta más rígido o menos rígido?

    \item Si dos resortes se colocan en serie, ¿cómo cambia la constante equivalente respecto de las constantes individuales?

    \item ¿Por qué un resorte real de masa apreciable no queda completamente descrito por
    \[ T=2\pi\sqrt{\frac{m}{k}}? \]

    \item ¿Qué observación experimental indicaría que un resorte está dejando de obedecer la ley de Hooke?
\end{enumerate}


\begin{resumen}
Un bloque de masa \(m\) unido a un resorte ideal de constante elástica \(k\) experimenta una fuerza restauradora:

\[ F_{s,x}=-kx. \]

Aplicando la segunda ley de Newton:

\[ ma_x=-kx. \]

Por tanto:

\[ \boxed{a_x=-\frac{k}{m}x}. \]

Comparando con la ecuación general del movimiento armónico simple:

\[ a_x=-\omega^2x, \]

se obtiene:

\[ \boxed{\omega=\sqrt{\frac{k}{m}}}. \]

El periodo y la frecuencia son:

\[ \boxed{T=2\pi\sqrt{\frac{m}{k}}}, \]

\[ \boxed{f=\frac{1}{2\pi}\sqrt{\frac{k}{m}}}. \]

Una masa mayor produce oscilaciones más lentas, mientras que un resorte más rígido produce oscilaciones más rápidas.

En un sistema vertical, la posición de equilibrio se encuentra mediante:

\[ \boxed{k\Delta\ell_{\mathrm{eq}}=mg}. \]

Al medir el desplazamiento desde esa posición, la dinámica vuelve a ser:

\[ ma_y=-ky. \]

Por ello, el periodo vertical ideal continúa siendo:

\[ \boxed{T=2\pi\sqrt{\frac{m}{k}}}. \]

Para varios resortes en paralelo:

\[ \boxed{k_{\mathrm{eq}}=\sum_i k_i}. \]

Para resortes en serie:

\[ \boxed{\frac{1}{k_{\mathrm{eq}}}=\sum_i\frac{1}{k_i}}. \]

En todos los casos, si la fuerza neta puede escribirse como:

\[ F=-k_{\mathrm{eq}}x, \]

el sistema realiza movimiento armónico simple con:

\[ \boxed{\omega=\sqrt{\frac{k_{\mathrm{eq}}}{m}}}. \]

Estas relaciones son válidas mientras el comportamiento sea aproximadamente lineal, las pérdidas puedan despreciarse y el modelo de resorte ideal sea adecuado.
\end{resumen}